\chapter{Ponte: testes em um projeto Quarkus}
\label{cap:ponte-quarkus}

Este capítulo liga o que vimos nas duas aulas ao Quarkus, que é o tema da
Parte~III. Só entra aqui o que foi confirmado na documentação oficial em
18/10/2026; o que não foi confirmado está marcado.

\section{\texttt{@QuarkusTest}}

Segundo o guia \citetitle{doc:quarkus-testing}, ao rodar um teste anotado
com \texttt{@QuarkusTest} a aplicação é iniciada antes do teste: o Quarkus
sobe uma vez antes do primeiro teste, todos os testes rodam, e ele é
desligado no fim. Os testes são \emph{CDI beans} completos, então é
possível injetar beans da aplicação neles.

O guia proíbe misturar \texttt{@QuarkusTest} com
\texttt{@QuarkusDevModeTest}, \texttt{@QuarkusProdModeTest} ou
\texttt{@QuarkusExtensionTest} na mesma execução e diz que, no momento,
\texttt{@QuarkusTest} e \texttt{@QuarkusIntegrationTest} não devem rodar na
mesma execução (em Maven, recomenda o Surefire para o primeiro e o Failsafe
para o segundo). O \texttt{@QuarkusIntegrationTest} serve para executar e
testar o artefato produzido pelo build, como caixa preta, e o guia diz que
ele tem os mesmos recursos e limitações do teste de executável nativo, cujo
guia exclui a injeção nos testes; o guia não afirma explicitamente que essa
exclusão vale para o modo JVM.

Na linguagem da aula 1: o \texttt{@QuarkusTest} é um teste de integração
com o SUT inteiro de pé; use-o nas bordas e mantenha as regras de negócio
em testes unitários rápidos.

\section{Dev Services e Testcontainers}

O guia \citetitle{doc:quarkus-dev-services} descreve o provisionamento
automático de serviços não configurados nos modos de desenvolvimento e de
teste. O Quarkus costuma iniciá-los usando o Testcontainers por baixo
(\citetitle{doc:testcontainers}). Eles sobem quando você inclui uma
extensão que os suporta e não configura conexão com um serviço externo.

Dois pontos práticos do guia: a maioria exige um ambiente de contêineres
(Docker ou Podman) e, para desligar todos os Dev Services, define-se
\texttt{quarkus.devservices\allowbreak.enabled=false}. O guia de testes acrescenta que
o \texttt{@QuarkusIntegrationTest} funciona com contêineres iniciados por
Dev Services.

\section{PIT em um projeto Quarkus}

A documentação do PIT mostra a configuração em Maven: o plugin
\texttt{pitest-maven} e o objetivo \texttt{mutationCoverage}, com
\texttt{targetClasses} para limitar os pacotes (\citetitle{doc:pitest-maven}).
\naoverificado{A integração específica do PIT com testes \texttt{@QuarkusTest}
não foi encontrada em documentação oficial do Quarkus ou do PIT; teste no
seu projeto antes de depender dela.}

\section{SmallRye Fault Tolerance e Resilience4j}

O Quarkus usa o SmallRye Fault Tolerance, que a documentação descreve como
uma implementação da especificação Eclipse MicroProfile Fault Tolerance
(\citetitle{doc:smallrye-fault-tolerance}), com uma API declarativa por
anotações e uma API programática. O guia
\citetitle{doc:quarkus-fault-tolerance} mostra \texttt{@Retry},
\texttt{@Timeout}, \texttt{@Fallback} e \texttt{@CircuitBreaker}, e descreve
o \emph{bulkhead} como o limite de execuções concorrentes.

O \citetitle{doc:resilience4j} é uma biblioteca de tolerância a falhas
pensada para programação funcional, com módulos de CircuitBreaker,
Bulkhead, Retry, RateLimiter e TimeLimiter, usados como decoradores de
lambdas e referências de método.

Os dois cobrem padrões de resiliência da aula 2 (retry, circuit breaker,
bulkhead). Nenhuma das duas documentações cita a outra,
\naoverificado{então não afirmamos equivalência um a um entre as
anotações do SmallRye e os módulos do Resilience4j.} O estilo é que
muda: anotações em métodos de beans CDI de um lado, decoradores
compostos no código do outro.

\section{Gatling contra uma aplicação Quarkus}

A documentação do Gatling descreve o SDK Java e, no modelo aberto,
\texttt{injectOpen}; as assertions quebram a simulação por percentil e por
taxa de falhas (\citetitle{doc:gatling-injection},
\citetitle{doc:gatling-assertions}). Um exemplo adaptado do formato mostrado na documentação (a página define \texttt{percentile4} como o percentil 99 por padrão, configurável em \texttt{gatling.conf}):

\begin{lstlisting}[language=Java]
setUp(scn.injectOpen(injectionProfile))
  .assertions(
    global().responseTime().percentile4().lt(1000),
    global().failedRequests().percent().lte(5.0)
  );
\end{lstlisting}

\naoverificado{Não encontramos guia oficial do Quarkus nem do Gatling sobre
usar um contra o outro.} Na prática, o cenário de carga é escrito como
em qualquer aplicação; a aula 2 recomenda rodar em homologação, com
limiares que quebram o build.
